\documentclass{article}
\usepackage{amsmath,amssymb}
\usepackage[vocab={numbers,specialnumbers}]{leanify}
\newtheorem{lemma}{Lemma}
\begin{document}

\section{Notation}

We write $F_n$ for the $n$th Fermat number.
We write $M_n$ for the $n$th Mersenne number.
We write $R_n$ for the $n$th repunit.
We write $P_n$ for the $n$th pentagonal number.
We write $T_n$ for the $n$th tetrahedral number.

\section{Kinds of integers}

\begin{lemma}\label{fermat}
$F_0=3$, $F_1=5$ and $F_2=17$. Every Fermat prime is odd.
\end{lemma}

\begin{lemma}\label{mersenne}
Let $n$ be a positive integer. If $M_n$ is prime, then $n$ is prime.
\end{lemma}

\begin{lemma}\label{twin}
$5$ is a twin prime, $11$ is a Sophie Germain prime and $23$ is a safe prime.
\end{lemma}

\begin{lemma}\label{squares}
$13$ is a sum of two squares, $64$ is a perfect cube, $27$ is a perfect power and $10$ is triangular.
\end{lemma}

\begin{lemma}\label{powers}
$16$ is a power of two and $81$ is a power of $3$. Every power of two $x>1$ is even.
\end{lemma}

\begin{lemma}\label{figurate}
$R_3=111$, $P_3=12$ and $T_3=10$. For every nonnegative integer $n$, $P_n\ge0$.
\end{lemma}

\begin{lemma}\label{fermat-two-squares}
Every prime $p\equiv1\pmod4$ is a sum of two squares.
\end{lemma}

\begin{lemma}\label{sums-of-squares}
Every prime $p$ with $p\equiv1\pmod 4$ is a sum of two squares. The number $7$ is not a sum of three squares. Every positive integer is a sum of four squares.
\end{lemma}

\begin{lemma}\label{cauchy-schwarz}
Let $n$ be a positive integer, let $x_1,\dots,x_n$ be real numbers and let $y_1,\dots,y_n$ be real numbers. Then $\left(\sum_{i=1}^n x_iy_i\right)^2\le\sum_{i=1}^n x_i^2\sum_{i=1}^n y_i^2$.
\end{lemma}

\end{document}
